%%% BEGIN fig:splitting-vertices
\begin{figure}[htbp]
\centering
\begin{tikzpicture}[fd]
  % ---------- (a) q -> q g ----------
  \begin{scope}
    \coordinate (V) at (0,0);
    \draw[quark] (-1.5,0) -- (V);
    \draw[quark] (V) -- (1.35,0.8);
    \draw[gluon] (V) -- (1.35,-0.8);
    \node[vertex] at (V) {};
    \node[left] at (-1.5,0) {$q$};
    \node[right] at (1.35,0.8) {$q\;\;(z)$};
    \node[right] at (1.35,-0.8) {$g\;\;(1-z)$};
    \node[font=\small,align=center,anchor=north] at (0.45,-1.25)
      {(a)\enspace$q\to qg$\\[2pt]$P_{qq}(z)=C_F\Big[\dfrac{1+z^2}{1-z}\Big]_+$};
  \end{scope}
  % ---------- (b) q -> g q ----------
  \begin{scope}[shift={(7.6,0)}]
    \coordinate (V) at (0,0);
    \draw[quark] (-1.5,0) -- (V);
    \draw[gluon] (V) -- (1.35,0.8);
    \draw[quark] (V) -- (1.35,-0.8);
    \node[vertex] at (V) {};
    \node[left] at (-1.5,0) {$q$};
    \node[right] at (1.35,0.8) {$g\;\;(z)$};
    \node[right] at (1.35,-0.8) {$q\;\;(1-z)$};
    \node[font=\small,align=center,anchor=north] at (0.45,-1.25)
      {(b)\enspace$q\to gq$\\[2pt]$P_{gq}(z)=C_F\,\dfrac{1+(1-z)^2}{z}$};
  \end{scope}
  % ---------- (c) g -> q qbar ----------
  \begin{scope}[shift={(0,-4.3)}]
    \coordinate (V) at (0,0);
    \draw[gluon] (-1.5,0) -- (V);
    \draw[quark] (V) -- (1.35,0.8);
    \draw[quark] (1.35,-0.8) -- (V);
    \node[vertex] at (V) {};
    \node[left] at (-1.5,0) {$g$};
    \node[right] at (1.35,0.8) {$q\;\;(z)$};
    \node[right] at (1.35,-0.8) {$\bar q\;\;(1-z)$};
    \node[font=\small,align=center,anchor=north] at (0.45,-1.25)
      {(c)\enspace$g\to q\bar q$\\[2pt]$P_{qg}(z)=T_F\big[z^2+(1-z)^2\big]$};
  \end{scope}
  % ---------- (d) g -> g g ----------
  \begin{scope}[shift={(7.6,-4.3)}]
    \coordinate (V) at (0,0);
    \draw[gluon] (-1.5,0) -- (V);
    \draw[gluon] (V) -- (1.35,0.8);
    \draw[gluon] (V) -- (1.35,-0.8);
    \node[vertex] at (V) {};
    \node[left] at (-1.5,0) {$g$};
    \node[right] at (1.35,0.8) {$g\;\;(z)$};
    \node[right] at (1.35,-0.8) {$g\;\;(1-z)$};
    \node[font=\small,align=center,anchor=north] at (0.45,-1.25)
      {(d)\enspace$g\to gg$\\[2pt]$\begin{aligned}P_{gg}(z)={}&2C_A\Big[\dfrac{z}{(1-z)_+}+\dfrac{1-z}{z}+z(1-z)\Big]\\[-2pt]&+\dfrac{\beta_0}{2}\,\delta(1-z)\end{aligned}$};
  \end{scope}
\end{tikzpicture}
\caption{领头阶分裂函数对应的四种 $1\to2$ 分裂。母部分子 $j$（左）分裂为携带其加号动量分数 $z$ 的部分子 $i$（右上）与携带 $1-z$ 的部分子（右下）；
$P_{ij}(z)$ 是 DGLAP 方程 $\partial f_i/\partial\ln\mu^2=\frac{\as}{2\pi}\sum_jP_{ij}\otimes f_j$ 中 $j\to i$ 的领头阶核。
$[\,\cdot\,]_+$ 为加号分布，$C_F[(1+z^2)/(1-z)]_+=C_F[(1+z^2)/(1-z)_++\tfrac32\delta(1-z)]$；$\beta_0=\tfrac{11}{3}C_A-\tfrac43T_Fn_f$。
(c) 中 $P_{qg}$ 对每一味夸克与反夸克分别计入；(c) 中反夸克线的箭头指向顶点（费米子数流向）。}
\label{fig:splitting-vertices}
\end{figure}
%%% END fig:splitting-vertices
